\chapter*{背景知识与学习路线}
\addcontentsline{toc}{chapter}{背景知识与学习路线}

学习本册不需要已经学过完整的计算机体系结构、操作系统、编译原理或高性能计算。但必须愿意把一个程序拆开来看：源码是一层，编译器是一层，汇编和 ABI 是一层，CPU 执行是一层，测量方法又是一层。本册要做的事，就是把这些层连成一条能推导、能观察、能复现实验的路径。

\section*{已经需要具备的基础}

开始学习本册前，最好已经具备下面这些 C/C++ 基础。它们不需要达到专家水平，但不能完全陌生：

\begin{itemize}
  \item 能写、编译、运行一个多文件 C/C++ 程序。
  \item 理解变量、表达式、函数、作用域、头文件和源文件。
  \item 理解数组、指针、引用、结构体、类和基本对象生命周期。
  \item 知道整数类型、浮点类型、\code{sizeof}、\code{alignof} 的基本用法。
  \item 能使用命令行运行程序、查看文件、重定向输出。
  \item 能读懂简单循环、条件分支、函数调用和基础标准库容器。
\end{itemize}

如果这些基础还不稳定，本册仍然可以读，但要放慢速度。每遇到一个 C/C++ 语法点，都要先写最小程序验证，而不是直接跳到汇编或性能结论。

\section*{本册要补齐的五类背景}

\subsection*{一、底层硬件背景}

硬件背景不是要求你设计 CPU，而是要求你知道程序最终运行在真实机器上。第一册会建立下面这些最小硬件模型：

\begin{itemize}
  \item bit、byte、word、地址、寄存器、内存和指令。
  \item CPU 是一个会读取指令、更新状态的机器，不是直接理解 C++ 的解释器。
  \item ISA 是软件和硬件之间的合同，微架构是实现这个合同的具体方式。
  \item load/store、分支、函数调用、栈和返回地址都是可观察的机器行为。
  \item cache、TLB、分支预测、流水线和乱序执行先作为性能直觉出现，深入模型留到第二册。
\end{itemize}

本册对硬件知识的要求是“能解释现象”，不是“能背参数”。例如，看到一个数组求和函数，应能说明它为什么需要地址计算、load、add、循环控制和返回值寄存器；至于某代 CPU 的端口分配和具体延迟，属于后续进阶内容。

\subsection*{二、软件系统背景}

软件背景的核心是：程序不是从源码直接运行。它会经过工具链和操作系统提供的多层机制：

\begin{itemize}
  \item 预处理、编译、优化、汇编、链接、加载和运行时启动。
  \item 翻译单元、目标文件、符号、重定位、可执行文件和共享库。
  \item 进程、虚拟地址空间、栈、堆、全局区、只读区和动态映射。
  \item \code{main} 不是进程第一条指令，运行时会在调用 \code{main} 前完成初始化工作。
  \item 编译器优化可能删除、移动、内联或改写代码，只要可观察行为符合语言规则。
\end{itemize}

这类知识对性能学习非常关键。很多错误 benchmark 不是算法慢，而是测到了初始化、动态链接、I/O、随机数生成、内存分配、页错误或被编译器删除后的空循环。

\subsection*{三、算法和数据背景}

本册不会讲复杂算法设计，但会训练你用底层视角重新看基础算法和数据：

\begin{itemize}
  \item 一个循环不仅是语法结构，也是控制流、归纳变量、地址计算和内存访问模式。
  \item 一个数组不仅是容器，也是连续对象序列和一组可推导的字节地址。
  \item 一个结构体不仅是字段集合，也是 offset、padding、alignment 和 ABI 约束。
  \item 一个函数不仅是代码块，也是调用约定、参数传递、返回值位置和寄存器保存规则。
  \item 一个性能数字不仅是结果，也是输入规模、分布、热路径、时间源、统计口径和环境共同产生的证据。
\end{itemize}

第一册的算法训练目标很朴素：能把 \code{sum}、\code{copy}、\code{dot}、条件计数、简单查表、结构体数组访问这类基础 kernel 解释清楚。后续卷中的 SIMD、GEMM、attention 和算子融合，都会建立在这些基础 kernel 的可推导模型上。

\subsection*{四、汇编和 ABI 背景}

汇编不是单独背指令表。它是连接 C/C++、机器码、寄存器、内存和 ABI 的中间层。本册学习汇编时只抓住几个稳定问题：

\begin{itemize}
  \item 参数从哪里进来，返回值从哪里出去。
  \item 每条指令读哪些寄存器或内存，写哪些寄存器或内存。
  \item 地址表达式如何由 base、index、scale、displacement 组成。
  \item 条件跳转依赖哪些 flags，\code{cmp} 和 \code{test} 为什么不保存普通结果。
  \item \code{call} 和 \code{ret} 如何使用栈保存和恢复控制流。
  \item caller-saved、callee-saved、stack alignment 和 red zone 这些 ABI 规则为什么存在。
\end{itemize}

学完第一册后，不需要立刻手写大型汇编库，但必须能读懂普通优化汇编，能从 C++ 函数签名推导 ABI 位置，能写小型手写汇编函数，并能用 \code{nm}、\code{readelf}、\code{objdump}、GDB 证明它真的按预期工作。

\subsection*{五、实验和测量背景}

性能学习必须从一开始就建立实验纪律。本册把 benchmark 当成受控实验，而不是随手计时：

\begin{itemize}
  \item 先证明正确性，再讨论速度。
  \item 明确 workload、kernel、timed region 和 validation 的边界。
  \item 保存原始样本，不只保存一个汇总数字。
  \item 报告编译器、编译参数、CPU、操作系统、输入规模和运行命令。
  \item 用反汇编确认测到的代码确实存在。
  \item 说明结论适用范围，不能从一个输入规模推断所有场景。
\end{itemize}

本册不要求每个实验都达到论文级严谨，但要求每个结论都有证据链。一个合格结论的形态应当是：

\begin{center}
\begin{minipage}{0.9\linewidth}
\begin{verbatim}
在这台机器、这个编译器、这些参数、这个输入规模下，
这个 kernel 的热路径生成了这些核心指令，
重复测量得到这些样本，
因此可以在这些限制内说它表现为某种趋势。
\end{verbatim}
\end{minipage}
\end{center}

\section*{第一册的知识链}

第一册的学习顺序不是随意安排的。每一部分都为下一部分提供必要背景：

\begin{center}
\begin{minipage}{0.92\linewidth}
\begin{verbatim}
第 1-5 章：
    建立源码、工具链、机器执行、数据表示和调试工具的地基。

第 6-10 章：
    用 x86-64 汇编和 System V AMD64 ABI 解释真实二进制接口。

第 11-12 章：
    把前面的源码、汇编、机器行为变成可复现的性能实验。
\end{verbatim}
\end{minipage}
\end{center}

不要跳过第一部分直接读汇编。没有源码到进程、CPU 状态、数据表示和工具链背景，汇编很容易变成指令速查。也不要跳过汇编直接做 benchmark。没有二进制审计能力，你无法判断自己测到的是源码里想象的代码，还是编译器改写后的另一段代码。

\section*{每个阶段的最低产出}

\subsection*{学完第一部分后}

你应该能做到：

\begin{itemize}
  \item 画出从 \code{.cpp} 到正在运行进程的粗粒度流程。
  \item 解释 CPU 执行的是机器码字节，不是 C++ 源码。
  \item 解释寄存器、内存、\register{rip}、\register{rsp} 和 flags 的作用。
  \item 从 \code{a[i]} 推导出字节地址公式。
  \item 用 \code{gdb}、\code{objdump}、\code{readelf}、\code{nm} 做基础观察。
\end{itemize}

\subsection*{学完第二部分后}

你应该能做到：

\begin{itemize}
  \item 读懂普通 x86-64 Intel 语法反汇编。
  \item 从 C++ 函数签名推导前几个参数和返回值的位置。
  \item 解释整数运算、地址计算、条件跳转、循环、函数调用和 switch 的常见汇编形态。
  \item 判断一个手写汇编函数是否破坏 ABI。
  \item 写小型 C++ 与汇编互操作实验，并用工具证明符号、重定位和调用路径。
\end{itemize}

\subsection*{学完第三部分后}

你应该能做到：

\begin{itemize}
  \item 把临时代码升级成有输入配置、重复测量、CSV 输出和环境记录的 benchmark。
  \item 区分 setup、warmup、timed region、validation 和 reporting。
  \item 用 reference 或不变量证明结果正确。
  \item 用反汇编和必要的 \code{perf stat} 尝试支撑性能解释。
  \item 写出别人能复现实验的报告。
\end{itemize}

\section*{从问题到证据的学习方法}

本册每个概念都应按同一种方法学习：

\begin{enumerate}
  \item 先问它为什么存在。没有这个机制会有什么问题。
  \item 再看它给程序员提供了什么抽象。
  \item 然后看底层大致如何实现或如何表现出来。
  \item 接着写最小代码，用工具观察证据。
  \item 最后说明它对正确性、性能和调试有什么影响。
\end{enumerate}

例如学习 ABI 时，不应只背“前六个整数参数放在 \register{rdi}、\register{rsi}、\register{rdx}、\register{rcx}、\register{r8}、\register{r9}”。还要知道为什么需要统一调用约定，调用者和被调用者如何分工，栈为什么要对齐，返回大结构体为什么可能出现隐藏指针，手写汇编破坏规则后会产生什么错误，以及如何用反汇编和 GDB 证明。

\section*{不追求一次理解所有细节}

第一册会反复预告 cache、TLB、分支预测、编译器优化、SIMD、PMU 和 AI 算子，但不会在这里彻底展开。预告的目的不是提前堆高级内容，而是说明：今天学的指针、对象布局、ABI 和 benchmark，以后都会在更复杂的性能问题里重复出现。

学习时应接受一个事实：第一遍读懂主线，第二遍跑通实验，第三遍写出解释，第四遍回头修正理解。底层知识不是靠记忆一次完成的，而是靠“代码、工具、证据、报告”反复闭环建立起来的。
